\documentclass[11pt]{article}
\usepackage[margin=1.1in]{geometry}
\usepackage{amsmath,amssymb,amsthm}
\usepackage[colorlinks=true,linkcolor=blue!50!black,citecolor=blue!50!black]{hyperref}

\newtheorem{theorem}{Theorem}[section]
\newtheorem{lemma}[theorem]{Lemma}
\newtheorem{proposition}[theorem]{Proposition}
\newtheorem{corollary}[theorem]{Corollary}
\theoremstyle{definition}
\newtheorem{definition}[theorem]{Definition}
\newtheorem{remark}[theorem]{Remark}
\newtheorem{obstruction}[theorem]{Obstruction}

\newcommand{\ch}{\operatorname{ch}}
\newcommand{\NS}{\operatorname{NS}}
\newcommand{\Pic}{\operatorname{Pic}}
\newcommand{\End}{\operatorname{End}}
\newcommand{\Q}{\mathbb{Q}}
\newcommand{\Z}{\mathbb{Z}}
\newcommand{\C}{\mathbb{C}}
\newcommand{\WQ}{W_{\Q}}
\newcommand{\vol}{\mathrm{vol}}
\newcommand{\FMc}{\widehat\Phi}

\title{The Decision Attempt in the Charge Frame:\\
Both Polarities at Their Certified Obstruction Points,\\
Sector Rigidity of the Fourier--Mukai Transform (Lemma W7),\\
and the Funnel to the Vertex\\
\large GST Hodge campaign, Cycle 4, Strike 2 (the decision attempt; route (b), ULTRA MODE)}
\date{2026-10-26}

\begin{document}
\maketitle

\begin{abstract}
\noindent
Cycle 4 strike 1 delivered the complete wall system on the pair channel
(Lemma W0, Theorems W4, W5, W6) and compressed the decision frame to a single
rigidity question with both polarities armed. This strike EXECUTES the decision
attempt in the charge frame and drives both polarities to their certified
obstruction points. On the NO route (charge rigidity) the six attack lines
(moduli-connectedness, corridor-invisibility kill chains, FM-closure via
product worlds, isogeny scaling, End-mass twisting, lattice/parity obstructions)
each die at a named committed wall, and the deaths are structural: every
additive invariant of a pure pair factors through the $(r,k,m,q)$-profile of
W4--W6, and that profile is consistent at every point of the certified
$(m,q)$-grid --- so no additive wall WITHIN the profile theory can
separate the charge strata (Theorem~\ref{thm:completeness}, the corridor
symmetry; the profile theory is the certified factorization of every
additive invariant through the $(r,k,m,q)$-data of W4--W6).
On the YES route the strike proves a new wall, \textbf{Lemma W7} (sector
rigidity of the cohomological Fourier--Mukai transform: the complement
structure preserves the NS/W/T sector decomposition; $\FMc(\theta^2)$ stays in
the NS-span, $\FMc(W_R)=\widehat W_R$, $\FMc(W_I)=\widehat W_I$), whose
corollaries kill the transform-transport and projector-correspondence
construction families: the algebraic projector $P_W$ cannot \emph{birth}
Weil content from NS/T content. The result is the \textbf{Funnel Theorem}
(\S\ref{sec:funnel}): within the committed wall system the dichotomy bit is
equivalent to the existence of a single W-charged algebraic surface class
(equivalently a holomorphic pure pair at some scale), and every mechanism
family in the committed arsenal generates $q=0$ content only. The decision
attempt therefore TERMINATES at the vertex: the campaign's certified state is
the reduction of the target on the open sector to one question --- does one
W-charged surface exist on the discriminator? Both resolution polarities stay
armed; neither fires inside the current system; the registry rows carry the
reopen conditions. No decision, no resolution, and no exhaustiveness claim
beyond the enumerated mechanism families is made (scope box, \S\ref{sec:scope}).
\end{abstract}

\section{Scope and discipline (read first)}\label{sec:scope}

\paragraph{What this dossier claims.}
\begin{enumerate}
\item \textbf{The NO-route obstruction map} (\S\ref{sec:no}): six attack lines
on charge rigidity, each driven to the certified wall where it dies.
\item \textbf{Lemma W7} (sector rigidity of $\FMc$; \S\ref{sec:w7}) with its
model certificate, and its two corollaries (the transform-transport kill; the
projector-birth kill).
\item \textbf{The additive-completeness theorem}
(Theorem~\ref{thm:completeness}): the W4--W6 profile is the COMPLETE additive
invariant theory of the pair channel, consistent at every certified grid
point; consequently any future deciding wall must be non-additive and
FM-closed.
\item \textbf{The Funnel Theorem} (\S\ref{sec:funnel}): the equivalence of the
bit with the W-charged-surface existence, plus the statement that every
enumerated mechanism family of the committed arsenal generates $q=0$ content.
\end{enumerate}

\paragraph{What this dossier does NOT claim.}
\begin{itemize}
\item Any decision of $S=0$ versus $S=W_\Q$ on any instance; any resolution of
the campaign target in either polarity.
\item Any exhaustiveness claim about conceivable pair constructions BEYOND the
enumerated mechanism families (the funnel covers the committed arsenal:
divisor intersections, tensor transports, FM transforms of NS-sector inputs,
$P_W$-correspondences, isogeny transport, Segre/pushforward on NS-sector
cycles).
\item Any claim that the two routes are closed permanently: the registry rows
(\S\ref{sec:terminal}) state the reopen conditions; a structurally new mechanism
reopens the corresponding route by the committed registry rules.
\item Any monoid-bit movement, any new seed, any change to the committed
equivalences.
\end{itemize}

\paragraph{Standing leaves.} All scheme-level inputs are the committed ones
(E3 GAGA, E5 Chern multiplicativity, E8 HRR/GRR, E9 Serre duality, E13 FM
toolkit, E14 Mukai band theorem, the dual-world axioms, the $K^0$ presentation,
the spectrum census as a record-level hypothesis in the projector idempotence,
the W2 pair-form equivalence content). Nothing certified in the green layers is
re-proved here; everything certified is cited with its receipt.

\section{Setting: the committed frame}\label{sec:setting}

The discriminator $A$ (nonsplit $\Q(i)$-Weil fourfold, $\NS(A)=\Z\theta$),
the model (D1)--(D3), the certified constants
($\int\theta^4=24$, $\int W_R^2=\int W_I^2=8$, $\int W_RW_I=0$,
$\theta^j\wedge W=0$ as classes for $j\ge 1$, $W\in\WQ$), the committed
equivalences (W1 dichotomy $S=0$ or $S=W_\Q$; W2 pair form $\iff$ algebraic;
on the discriminator pair form $\iff S=W_\Q$; W5 the bit is duality-invariant),
and the strike-1 walls on the pair channel (W4 invisibility corridor: the
difference-Euler of a pure pair vanishes against every line bundle, the only
linear detector being $m\langle W_R,\ch_2(G)\rangle$; W6 the pair Chern table:
equal $(r,c_1,c_3)$, $c_2$-difference $-mW_R$, charge split
$q(E_+)-q(E_-)=-8m$, the mixed-End table). The Lean layer
(\texttt{GSTPairChannelWallsCertificate.lean}) kernel-verifies the arithmetic
of all of these in the bitmask exterior model; its CI arc is recorded in
\S\ref{sec:verif}.

\section{The NO route: charge rigidity --- the six-line obstruction map}\label{sec:no}

The negative certificate requires: \emph{for every two holomorphic bundles with
shared $(r,c_1\in\Z\theta,c_3)$ and the $c_4$-relation of W6(iii), the Weil
charges agree} (then no pure pair exists at any scale, so $S=0$ by the
committed W2(iii), and $W_R,W_I$ are nonalgebraic by W1(iii) --- the PROMPT
\S2 negative resolution). Six attack lines, six certified deaths:

\begin{obstruction}[Moduli-connectedness]
Chern classes are locally constant in flat families over connected bases, so
$q$ is a locally constant function on the bundle moduli. A charge split $-8m$
is a jump between components --- permitted. Rigidity would require the moduli
to be connected across the charge strata or the strata to be empty; neither is
reachable from the committed walls. \emph{Dies at: nothing available --- the
needed statement is precisely the unproved rigidity itself.}
\end{obstruction}

\begin{obstruction}[Corridor-invisibility kill chains]
The carrier-layer kill (V2: single band $+$ Euler-zero $\Rightarrow$
acyclic-for-all-$P$ $\Rightarrow$ $\FMc(F)=0$ $\Rightarrow$ $F=0$) needs a
\emph{bundle} with a WIT band structure. The pair difference
$x=[E_+]-[E_-]\in K^0$ carries no band structure, and W4 makes every additive
probe blind to it at every scale and every polarization level.
\emph{Dies at: W4 --- the corridor is exactly the region where the kill chain
does not reach.}
\end{obstruction}

\begin{obstruction}[FM-closure via product worlds]
Testing the pair on a product world ($A\times\widehat A$, or split Weil worlds)
exits the corridor: at $\NS$-rank $\ge 2$ the annihilation Lemma W0 fails
(the committed consistency probe), and the walls are corridor-local. Every
attempt to manufacture a contradiction off-world loses the walls.
\emph{Dies at: W0's own failure locus --- the corridor has no exit.}
\end{obstruction}

\begin{obstruction}[Isogeny scaling]
Pairs transport along isogenies with the scale multiplied by the Weil-pullback
factor; with W5 (duality transports at the SAME scale) and the committed
principality (Theorem M: one principal scale divides every scale), the bookkeeping
is consistent at every level; no divisibility contradiction is producible.
\emph{Dies at: T1--T5 $+$ M --- all consistent.}
\end{obstruction}

\begin{obstruction}[End-mass twisting]
W6(v) expresses every End-Euler difference as a linear form
$8m^2-m\,q(E_-)$, $-m\,q(E_-)$, $8m^2-2m\,q(E_-)$ in the FREE integer
$q(E_-)\in\Z$; the synthetic instances realize the table at $q$-values $0$ and
$\pm 8$, and the $(m,q)$-grid certificate (\S\ref{sec:verif}) realizes it at
every grid point. No sign, parity, or positivity obstruction exists: the
carrier layer is the $q=0$ slice and every other $q$ is a priori live.
\emph{Dies at: W6(v) --- $q$ is free.}
\end{obstruction}

\begin{obstruction}[Lattice and parity (the Gr(2,5) pattern)]
The Grassmannian determination ($m^*=360$) came from a $K$-lattice necessity.
On the discriminator the committed minimal scale is $m_0=1$ (K\"unneth
unimodularity of the torus $K$-lattice, committed strike-1): the pair table is
integral-feasible at $m=1$, and the $P^2$-parity mechanism has no analogue here
(the certified $P^2$ device lives on the line-sum instance, not on the fourfold
lattice). \emph{Dies at: the committed $m_0=1$.}
\end{obstruction}

\begin{theorem}[Additive completeness; the corridor symmetry]\label{thm:completeness}
Let $x=[E_+]-[E_-]$ be a pure pair for $W_R$ at scale $m\ge 1$ on the
discriminator. Then:
\begin{enumerate}
\item[(i)] every additive invariant $\chi(x\otimes G)$ ($G$ perfect) equals
$m\langle W_R,\ch_2(G)\rangle$ --- the W4(iii) detector is the COMPLETE linear
theory;
\item[(ii)] the full additive profile of the pair is the quadruple
$(r,k,m,q)$ with the W6 relations (shared $(r,c_1,c_3)$, $c_2$-split $-mW_R$,
charge split $-8m$, mixed-End table), and the profile is consistent at every
point of the certified grid $m\in\{1,2,3,4\}$, $q/8\in\{0,1,2,3\}$,
$c_1\in\{0,\theta,2\theta\}$, $\gamma\in\{0,\vol\}$ (96 instances, zero
failures);
\item[(iii)] consequently no additive wall \emph{within the profile theory}
--- no necessary condition expressed through the twisted Euler
characteristics that the profile computes --- can separate the charge
strata or exclude the pair: separation would require a profile value to be
unrealizable, and the profile is internally consistent (the universal laws
admit the table, and the grid certificate realizes it throughout the
certified range $m\in\{1..4\}$, $q/8\in\{0..3\}$). Any future deciding
wall must therefore operate OUTSIDE the profile theory (be non-additive,
i.e.\ constrain data the profile does not compute), and by W5 it must in
addition be FM-closed. This is a statement about the certified profile
theory, not an exhaustiveness claim about all conceivable additive
constraints (scope box).
\end{enumerate}
\end{theorem}

\begin{proof}
(i) is W4(iii) with the W0-annihilated cross-grades. (ii): the universal laws
of the table are the hypothesis-form theorems of the Lean layer
(charge split from $\int W_R^2=8$; the mixed-End lines from the split and the
annihilated $c_1^2$-pairing; the third line as the telescoped sum of the first
two), and the grid certificate realizes every profile point explicitly with
grade-by-grade purity. (iii): a wall within the profile theory is a statement about
the values $\chi(x\otimes G)$; by (i)--(ii) these values are determined by
$(r,k,m,q)$ through the table relations (universally, as hypothesis-form
laws); the determined values are internally consistent -- the universal laws
admit the table by construction, and the grid certificate realizes every
profile point in the certified range with grade-by-grade purity; hence no
constraint expressible within the profile can exclude the pair. A constraint
that excludes the pair must therefore constrain data OUTSIDE the profile
(non-additive data), which is the stated conclusion. No claim is made about
the nonexistence of non-profile additive-flavored constraints (that would be
an exhaustiveness claim; the scope box excludes it).
\end{proof}

\begin{remark}
Theorem~\ref{thm:completeness} is the precise form of the statement ``the NO
route is unreachable within the certified profile theory'': it does not say
rigidity is false --- it says no constraint the profile theory can express
can establish it.
\end{remark}

\section{Lemma W7: sector rigidity of the Fourier--Mukai transform}\label{sec:w7}

The degree-$4$ rational cohomology of the Weil world decomposes into the three
sectors
\[
H^4(A,\Q)\;=\;\underbrace{\Q\theta^2}_{\mathrm{NS\ image}}\;\oplus\;
\underbrace{\WQ}_{\mathrm{Weil}}\;\oplus\;T,
\]
where $T$ is the complementary (transcendental) sector. In the exterior model
the cohomological FM transform $\FMc$ is the wsign-weighted complement on
degree-$4$ monomials (committed W5 normalization).

\begin{lemma}[Sector rigidity]\label{lem:w7}
The complement structure preserves the model sectors: the complement of a
block-union word is a block-union word ($\mathrm{NS}\to\mathrm{NS}$); the
complement of a one-letter-per-block word is a one-letter-per-block word of
the same parity ($W\to W$); the complement of a $T$-word is a $T$-word.
Consequently $\FMc$ preserves the sector decomposition:
\[
\FMc(\Q\theta^2)=\Q\widehat\theta^2,\qquad
\FMc(\WQ)=\widehat\WQ,\qquad \FMc(T)=\widehat T,
\]
and in particular $\FMc(\theta^2)$ lies in the dual NS-span,
$\FMc(W_R)=\widehat W_R$ and $\FMc(W_I)=\widehat W_I$ exactly.
\end{lemma}

\begin{proof}
The census: $70$ degree-$4$ words $=$ $6$ block-union words $+$ $16$ Weil
words $+$ $48$ $T$-words. Complementing a block-union word: full blocks and
empty blocks swap, so the complement of $\{a,b\}$-full is
$\{c,d\}$-full --- again a block-union. Complementing a one-letter-per-block
word flips the letter in every block, so the number of $f$-letters goes to
$4-\#f$: the parity is preserved, the even words (the $W_R$ monomials)
permute among themselves and the odd words (the $W_I$ monomials) among
themselves. Since the complement map is an involution and the two model
subspaces are closed, the set-theoretic complement of their union ($T$) is
closed as well. Linearity gives the sector statements; the class-level
identities $\FMc(W_R)=\widehat W_R$, $\FMc(W_I)=\widehat W_I$ are the committed
W5 normalization (kernel-certified in the Lean layer), and
$\FMc(\theta^2)\in\Q\widehat\theta^2$ is verified exactly in the model
certificate (the transform of $\theta^2$ is $\widehat\theta^2$ with the
expected coefficients). The scheme-level identification of the model sectors
with $\Q\theta^2\oplus\WQ\oplus T$ runs through the committed leaves as always.
\end{proof}

\begin{corollary}[Transform-transport kill]\label{cor:w71}
The FM orbit of the NS-sector never meets the Weil sector: no Fourier--Mukai
construction whose input has zero Weil content (in particular every
line-bundle input, every tensor transport of NS-sector classes, and every
FM-dual of such) produces Weil content. The transform-transport family of
conceivable YES-constructions is dead.
\end{corollary}

\begin{corollary}[Projector-birth kill]\label{cor:w72}
Assume the committed spectrum census (the record-level hypothesis of the
projector idempotence) and the correspondence-action compatibility (leaf E1).
Then for every $x\in K^0_{\mathrm{alg}}(A)$:
\[
\ch\bigl(P_W(x)\bigr)\;=\;P_W\bigl(\ch(x)\bigr)\;=\;\text{the Weil component
of }\ch(x),
\]
because the committed kernel-green spectral facts give $P_W=1$ on the Weil
plane, $P_W=0$ on the mixed factor, and $P_W(25)=0$ on the ambient
(NS-generated) factor, and the census extends the vanishing to the $T$-blocks.
Consequently the algebraic projector cannot \emph{birth} Weil content from
NS/T content: $P_W(x)$ has nonzero Weil character iff $\ch(x)$ already does.
The projector-correspondence family of conceivable YES-constructions is dead.
\end{corollary}

\begin{remark}
Corollary~\ref{cor:w72} sharpens the committed W1(iv) one-source machinery from
the opposite side: the projector \emph{reconstructs} explicit cycles from a
detected source, but it never \emph{creates} a source.
\end{remark}

\section{The YES route: the Funnel Theorem}\label{sec:funnel}

\begin{theorem}[Funnel]\label{thm:funnel}
On the discriminator, the following are equivalent:
\begin{enumerate}
\item[(i)] the dichotomy bit fires: $S=W_\Q$ (equivalently: $W_R$ is
algebraic);
\item[(ii)] there exists a W-charged algebraic surface class: a codimension-$2$
cycle $Z$ with nonzero Weil component of $\cl(Z)$;
\item[(iii)] there exists a holomorphic pure pair at some scale $m\ge 1$
(the committed W2 pair-form equivalence, with the W6 table as its
Chern-class face).
\end{enumerate}
Moreover, every mechanism family in the committed arsenal produces $q=0$
(zero-Weil-content) classes:
\begin{itemize}
\item divisor intersections: the NS-subalgebra is annihilated by $W$
(Lemma W0 --- the whole NS-subalgebra is blind to the Weil grades);
\item tensor/line-bundle transports: twisting by $L=\mathcal O(k\Theta)$
preserves the scale and adds only NS-sector content (W0 $+$ the twist check:
a pure pair twisted by any $L$ is a pure pair at the same scale, so no
twist ever creates Weil content);
\item Fourier--Mukai transforms: Lemma W7 / Corollary~\ref{cor:w71};
\item algebraic projector correspondences: Corollary~\ref{cor:w72};
\item isogeny transport: the committed T-laws transport the bit within an
isogeny class (scales multiply by the Weil-pullback factor); they cannot
transport it ACROSS non-isogenous classes, and the split worlds where
Weil's construction produces pairs are not isogenous to the discriminator
(committed);
\item Segre/pushforward constructions on NS-sector cycles: these realize
carriers at the $q=0$ slice (the committed carrier instance), and their
inputs (NS-sector cycle classes) carry no Weil component by W0.
\end{itemize}
Hence the entire remaining content of the YES route is the vertex: exhibit
one W-charged surface (equivalently one pure pair at any scale).
\end{theorem}

\begin{proof}
The equivalences (i)--(iii) are the committed W2(iii)/W1 with the W6 table
(attributed and recorded as normalization per the M3 rule; no new equivalence
content is claimed here). The mechanism list: each entry cites the wall that
kills it, as proven in the cited sections and the committed receipts. The
funnel statement is the conjunction: the bit is equivalent to the
W-charged-surface existence, and no enumerated mechanism reaches it.
\end{proof}

\begin{remark}[What would reopen the routes]
By the registry rules (\S\ref{sec:terminal}): the NO route reopens on a
structurally new mechanism in the sense of the committed registry standard --- e.g.
a non-additive, FM-closed rigidity invariant separating the charge strata
(refined Brill--Noether type data, Picard-sheaf monodromy invariants, or a
determinant-of-cohomology rigidity theory). The YES route reopens on any
W-charged surface construction: per-world vertex certificates (the route (a)
pattern of cycle $3$: lattice envelope $+$ witness production), deformations
across the Weil family that pin algebraicity at special members and transport
it, or non-abelian correspondence constructions beyond the committed
$f$-algebra. Both rows stay armed; neither polarity fired inside the current
system.
\end{remark}

\section{The certified terminal state of the decision attempt}\label{sec:terminal}

\begin{itemize}
\item The decision attempt EXECUTED in the charge frame with both polarities
at full force and duality armed (W5).
\item The NO route: unreachable within the certified profile theory
(additive $+$ transport $+$ lattice arsenal;
Theorem~\ref{thm:completeness}, with its scope statement); registry status
\textbf{BLOCKED}, reopen condition as above.
\item The YES route: funneled to the vertex (Theorem~\ref{thm:funnel}); registry
status \textbf{BLOCKED-funnel}, reopen condition as above.
\item The campaign's certified state is therefore the \textbf{reduction
stack}: HC on the open sector $\iff$ FS$_\forall$ (cycle 3, CI green, audit
closed) $\iff$ at the forcing vertex: the existence of one W-charged surface on
the discriminator (equivalently one pure pair at some scale, equivalently the
bit $S=W_\Q$).
\item Per the campaign contract (PROMPT \S2), the affirmative resolution is
the full quantified theorem and the negative resolution is the zero-test
certificate; the decision attempt shows both remain armed and neither is
reachable within the current system --- the funnel theorem is the terminal
certified statement of this strike, and the vertex is the campaign's
remaining target.
\end{itemize}

\section{Verification layer}\label{sec:verif}

\paragraph{The decision-attempt certificate script.}
\texttt{research/experiments/task14c\_decision\_attempt\_checks.py} (committed
alongside this dossier). Bounds: pure Python, exact \texttt{Fraction}
arithmetic, $\le 2^{10}$ monomial operations per instance, $96$ grid instances
$+$ the $70$-word census $+$ the sector-closure checks $+$ $4$
twist-preservation checks; runtime $<5$\,s, memory $<20$\,MB; no floats. It
verifies: (1) the W7 census ($6/16/48$) and all complement closures
(NS-union $\to$ NS-union; even Weil $\to$ even Weil; odd Weil $\to$ odd Weil;
$T\to T$), the class-level $\FMc(\theta^2)\in$ NS-span, and the committed
$\FMc(W_R)=W_R$, $\FMc(W_I)=W_I$; (2) the $(m,q)$-grid: purity
grade-by-grade, the charge split $-8m$, and the three mixed-End identities at
all $96$ points ($m\in\{1,2,3,4\}$, $c_2(E_-)=aW_R$ with $a\in\{0..3\}$,
$c_1\in\{0,\theta,2\theta\}$, $\gamma\in\{0,\vol\}$) --- zero failures;
(3) the twist-preservation of the pair table under
$L=\mathcal O(k\Theta)$, $k\in\{1,2\}$.

\paragraph{The Lean layer (the wall arithmetic).}
\texttt{GSTPairChannelWallsCertificate.lean}: the W0 class annihilations, the
period constants (second encoding), the W4 corridor arithmetic, the W5 FM
normalization $+$ round trips, and the W6 table (universal laws $+$ synthetic
instances $+$ carrier consistency $+$ detector instance), with the
pre-verification mirror
\texttt{research/experiments/task14a\_lean\_data\_mirror.py} (all DATA
cross-checked against an independent exterior reconstruction). CI arc: round 1
(run 40138713492) failed with three diagnosed causes --- (a) kernel-stuck
decides from non-proven primitives (bitwise ops and \texttt{flatMap}; replaced
by the committed div/mod $+$ structural-recursion idioms), (b) a sign slip in
the hypothesis-form third mixed-End line ($d=t_1-t_2$ corrected to the
telescoped $d=t_1+t_2$; the direct instance computations were correct
throughout; the kernel caught the packaging error), (c) a capstone projection
off by one conjunction level. Round 2 pushed with all three repairs; the green
receipt is appended to the record and INDEX when the run completes.

\paragraph{Attribution discipline (M3).} The funnel theorem's equivalence
content is the committed W2/W1 equivalence, attributed and recorded as
normalization; the NEW content of this strike is the obstruction map
(\S\ref{sec:no}), Lemma W7 with its corollaries, the additive-completeness
theorem, and the mechanism-list of the funnel. No equivalence is claimed as
new; no decision is claimed in either polarity.

\section{Provenance}

\begin{itemize}
\item Cycle 4, Strike 2 of the GST Hodge campaign (Mode B, V6 protocol; boss
directive 2026-10-26: route (b), ULTRA MODE, finish the work in this round).
\item Builds on: the committed Weil-forcing frame W1--W3 (cycle 2 strike 1,
audit-closed 9-c), the carrier-mechanism verdict V1--V4 (cycle 2 strike 2,
audit-closed 10-c), the for-all-N reduction (cycle 3, audit-closed 11-c), and
the pair-channel walls W0/W4/W5/W6 (cycle 4 strike 1, committed 147cfc4;
Lean layer in the CI arc recorded above).
\item New labels: Lemma W7 (sector rigidity), Corollaries W7.1/W7.2,
Theorems of \S\ref{sec:no}/\S\ref{sec:funnel} (the obstruction map, the
additive-completeness theorem, the funnel theorem). The label series extends
the committed W-frame; no committed label is reused.
\item No mathematical content of the committed state is altered; the research
files in this commit do not fire any workflow.
\end{itemize}

\end{document}
